% --- Archon physics formalization source begin ---
% archon:physics
% archon:covers IPhO2026Problems/problem_ipho_2026_t3_c2.lean
% archon:source-report reports/ipho_2026/problem_ipho_2026_t3_c2.source.json
% archon:problem-id ipho_2026_t3
% archon:part-id T3-C2
% archon:previous-part ipho_2026_t3_b1
% archon:previous-part-policy derive_inline_from_problem_only_material
% archon:previous-part ipho_2026_t3_c1
% archon:previous-part-policy derive_inline_from_problem_only_material

\chapter{Ipho Physics problem ipho\_2026\_t3\_c2}
\label{ch:IPhO2026Problems_problem_ipho_2026_t3_c2}

\paragraph{Problem source.}
\#\# Physical scenario

The paramagnetic torus executes the Carnot refrigeration cycle
1 -> 2 -> 3 -> 4 -> 1 shown in Figure 3b in the H-versus-T plane.  T\_h and T\_c
are the hot- and cold-reservoir temperatures; Q\_h is the magnitude of heat
delivered to the hot reservoir and Q\_c is the magnitude absorbed from the cold
reservoir.  The equation of state is T*M*V = n*K*H and the isothermal heat
relation from part B may be reused.

\#\# Current subquestion T3-C2

Express M\_1 in terms of M\_2, M\_3, and M\_4.

\paragraph{Shared problem context.}
The paramagnetic torus executes the Carnot refrigeration cycle
1 -> 2 -> 3 -> 4 -> 1 shown in Figure 3b in the H-versus-T plane.  T\_h and T\_c
are the hot- and cold-reservoir temperatures; Q\_h is the magnitude of heat
delivered to the hot reservoir and Q\_c is the magnitude absorbed from the cold
reservoir.  The equation of state is T*M*V = n*K*H and the isothermal heat
relation from part B may be reused.

\paragraph{Current subquestion.}
Express M\_1 in terms of M\_2, M\_3, and M\_4.

\paragraph{Figure/image path.}
ipho\_2026\_source/image/T3\_page-3.png

\paragraph{Reusable previous-part conclusions.}
\begin{itemize}
\item Source T3-B1. Question: At fixed temperature T, H changes from H\_i to H\_f. Find the heat Q transferred into the torus. Policy: derive\_inline\_from\_problem\_only\_material
\item Source T3-C1. Question: Label T\_h and T\_c on Figure 3b and identify the processes on which Q\_h and Q\_c are transferred. Policy: derive\_inline\_from\_problem\_only\_material
\end{itemize}

\paragraph{Formalization target.}
create a compiling Lean file with sorry bodies at `IPhO2026Problems/problem\_ipho\_2026\_t3\_c2.lean`.
The Lean declarations must preserve the physical quantities, dimensions or dimensional roles, figure labels, governing-law hypotheses, and final relation expressed by this problem.
Use Mathlib/Physlib names found through LeanExplore where available. If a domain API is missing, introduce faithful local abstractions rather than scalar placeholder aliases.
This is an answer-blind solve-phase task. Derive a candidate result only from the problem-side evidence above; do not assume a target value, select a tolerance around a desired value, or encode a desired result into a candidate domain.
For numerical questions, define the raw end-to-end quantity and apply an explicit source-derived rounding rule. For identification questions, characterize or prove uniqueness before recording the derived candidate.

\begin{theorem}[Ipho Physics formalization target]
\label{thm:physics:ipho_2026_t3_c2:target}
\leanok
The assigned autoformalize agent should translate this IPhO physics problem into Lean declarations in the covered file, with theorem and lemma proof bodies written as `by sorry`.
\end{theorem}
\begin{proof}
\leanok
Formalization completed in \texttt{IPhO2026Problems/problem\_ipho\_2026\_t3\_c2.lean}
(namespace \texttt{IPhO2026.T3C2}) and all proof bodies are now kernel-checked
(no \texttt{sorry}): the bridge lemmas
\texttt{heat\_cold\_leg\_magnetization\_form},
\texttt{heat\_hot\_leg\_magnetization\_form} (EOS solved for $H_i$ via
\texttt{eq\_div\_iff}, substituted into the T3-B1 heat relation, cleared by
\texttt{field\_simp}), \texttt{magnetization\_balance\_of\_heats}
(cross-multiplied Clausius relation, cancellation of the common prefactor
$A\,T_h\,T_c \neq 0$ by \texttt{mul\_left\_cancel₀}), the orientation bridges
\texttt{cold\_leg\_absorbs\_heat} / \texttt{hot\_leg\_rejects\_heat}
(\texttt{pow\_lt\_pow\_left₀}, positivity of the prefactor), and the main
targets \texttt{magnetization\_vertex\_one\_sq} and
\texttt{magnetization\_vertex\_one} are all proved, giving
$M_1^2 = M_2^2 - M_3^2 + M_4^2$ and $M_1 = \sqrt{M_2^2 - M_3^2 + M_4^2}$.
Axiom audit: only \texttt{propext}, \texttt{Classical.choice},
\texttt{Quot.sound}.
\end{proof}
% --- Archon physics formalization source end ---
